\subsection{A nonfree coefficient-tail Koszul closure}

Let $R$ be $\mathbb Z$ or $\mathbb Z_{(p)}$, and
$S_R=\Gamma_R(T)$. Fix $N\geq1$ and let $I_N$ be
its coefficient ideal of weights at least $N$.
For $m\geq1$ let $D_m\subset S_R$ be the ideal
generated by all products
\[
 \gamma_{a_1}(b_1)\cdots\gamma_{a_l}(b_l),
 \qquad b_i\in I_N,\quad a_i\geq1,
 \quad a_1+\cdots+a_l\geq m.
 \tag{AT.20}
\]
Put $D_0=S_R$. This is the usual divided-power
filtration of the specified ideal; (AT.20) is the
definition used here.

\begin{theorem}[AT10: Proved; a coefficient-tail closure with actual evaluation maps]
\label{AT10}
Let $M$ be a free $R$-module, $r\geq1$, and $b\geq0$.
Give two copies $M_x,M_y$ form degrees $2r-1,2r$
and weight $b$. Form the ordinary polynomial Koszul
algebra
\[
 F(M)=\Lambda_R(M_x)\otimes_R\operatorname{Sym}_R(M_y),
 \qquad dx(v)=y(v),\quad dy(v)=0.
\]
Write $F_m(M)$ for its total arity $m$ part.
Then
\[
 E_{N,r,b}(M)
   =\bigoplus_{m\geq0}D_m\otimes_R F_m(M)
       \ \subset\ S_R\otimes_R F(M)
 \tag{AT.21}
\]
is a strict PD differential algebra, with its
positive-weight ideal distinguished. Its terms
are free over $R$, and its coefficient inclusion
is a weight-preserving cochain homotopy equivalence.
It has support $q\leq w$ whenever $N+b\geq2r$.

Given any strict PD differential $S_R$-algebra $A$
and a weight-preserving $R$-linear map
$\alpha:M\to A^{b,2r-1}$, the ordinary evaluation
$x\mapsto\alpha$, $y\mapsto d\alpha$ restricts to
an actual strict PD DGA map
\[
 E_{N,r,b}(M)\longrightarrow A.
 \tag{AT.22}
\]
These maps are natural in $M$ and therefore are
simplicial when the parameters and $\alpha$ are
simplicial. No primitive is chosen in the target.

For a split inclusion of free modules $M\hookrightarrow M'$
with free complement, the map
$E_{N,r,b}(M)\to E_{N,r,b}(M')$ has a deformation
retraction whose homotopy is linear over the old
DG algebra. Consequently, for every strict PD map
$E_{N,r,b}(M)\to X$, the ordinary PD pushout
\[
 X\longrightarrow X\otimes_{E_{N,r,b}(M)}E_{N,r,b}(M')
 \tag{AT.23}
\]
is a weight-preserving cochain homotopy equivalence.
For degreewise split simplicial inclusions this
conclusion holds in each simplex. No termwise
flatness assertion for the pushout is included.
\end{theorem}

\begin{proof}
The definition gives $D_mD_l\subset D_{m+l}$.
It also gives $\gamma_a(D_m)\subset D_{am}$.
For completeness, on one product in (AT.20), apply
the scalar PD identity to one factor and the
iteration identity to its divided power. Each
original divided-power index is multiplied by $a$;
the coefficient is an integer. Multiplication by
an arbitrary coefficient from $S_R$ raises no
required index, and the scalar identity raises
that coefficient to the $a$th power. The PD sum
formula treats a sum of such products: a summand
with indices $a_1+\cdots+a_l=a$ belongs to
$D_{ma_1}\cdots D_{ma_l}\subset D_{ma}$.
This proves the asserted filtration compatibility.
Its generators all have coefficient weight at
least $mN$.

Embed (AT.21) in the rational graded-commutative
algebra obtained from $S_R\otimes F(M)$.
For a homogeneous even polynomial/exterior
monomial $v$ of arity $m>0$ and $a\in D_m$,
its rational divided power is the integral
expression
\[
 \gamma_n(av)=\gamma_n(a)v^n\in D_{mn}F_{mn}(M).
 \tag{AT.24}
\]
If $v$ contains exterior factors and has even
positive exterior degree, its powers of index
at least two are zero, consistent with this
formula. On the coefficient part use the given
PD structure of $S_R$. The sum formula now proves
that every even element of the positive-weight
ideal has all its divided powers in (AT.21).
All PD identities, the strict odd-product identity,
and differential compatibility hold rationally;
the displayed module is torsion free over $R$,
so they hold integrally. The differential preserves
arity and the coefficient ideals, proving the
strict PD DGA assertion. This construction is
independent of a basis because the rational
operations and the functorial submodule (AT.21)
are so.

To contract the positive arities, choose a basis
of $M$ and decompose by the arity in each variable.
For each positive multi-arity, take the first
variable with positive arity. Its polynomial
Koszul complex is the identity pair
\[
 x_i y_i^{a-1}\longmapsto y_i^a \qquad(a\geq1).
\]
Sending $y_i^a$ to $x_i y_i^{a-1}$ and the latter
to zero, with the tensor sign convention, contracts
that multi-arity component. The homotopy preserves
the total arity and hence preserves the coefficient
factor $D_m$ in (AT.21). It is $S_R$-linear and
weight preserving. It proves the coefficient
homotopy equivalence in every weight. Each weight
piece of each $D_m$ is a submodule of the rank-one
free $R$-module $S_R^w$, hence free over the PID
$R$; the monomial decomposition proves freeness
of the terms. An arity-$m$ term has weight at
least $m(N+b)$ and form degree at most $2rm$,
proving the support statement.

Ordinary evaluation on $S_R\otimes F(M)$ is a
well-defined DGA map: odd images have square zero,
the even images commute, and their differentials
are the prescribed ones. To check the PD identity
on its restriction, apply the target scalar PD
identity to (AT.24). The coefficient $a$ has
positive weight and belongs to the given coefficient
PD ideal, and the evaluated monomial is even.
Thus its divided power is exactly the evaluation
of $\gamma_n(a)v^n$. Coefficients and even sums
are treated by their defining PD identities.
This proves (AT.22), including targets with
arbitrary torsion. It also proves naturality,
since no chosen basis or primitive enters the
evaluation formula.

For the relative assertion, write $M'=M\oplus V$.
Projection to $M$ induces an algebra retraction
$\pi:E(M')\to E(M)$. On a component with positive
arity in $V$, use the same polynomial Koszul
contraction, choosing the first active variable
of $V$. Act as the identity on the old $F(M)$
factor and use the tensor sign rule. This homotopy
$H$ preserves total arity and satisfies
\[
 dH+Hd=1-i\pi,\qquad
 H(ae)=(-1)^{|a|}aH(e)\quad(a\in E(M)).
 \tag{AT.25}
\]
The second equation follows because the first
active new variable is unchanged by multiplying
by an old-variable monomial. The coefficient
filtration does not change this identity:
multiplying its coefficients uses precisely
$D_lD_m\subset D_{l+m}$. On the part with no
new variable set $H=0$. This proves all of
(AT.25) inside the specified modules.

We also justify the use of the ordinary tensor
in (AT.23). Graded augmented PD algebras have
presentations by free divided-power/exterior
algebras and PD-stable relation ideals. The
tensor product of two free presentations is
the free presentation on the disjoint generator
sets. Each original relation ideal remains
PD stable after adjoining the other factors,
by the scalar and sum identities. To impose
the common base, adjoin every relation
$f(c)\otimes1-1\otimes g(c)$. For a positive
even $c$, its divided power is zero modulo
these relations: in the PD sum expansion
replace each $f(\gamma_i c)$ by
$g(\gamma_i c)$, which is permitted because
the relation is imposed for every $c$, including
all its divided powers. The resulting expansion
is $\gamma_n(g(c)-g(c))=0$. Products involving
odd relation generators use the strict
odd-product identity; even weight-zero
relations use the scalar identity. The full
relation ideal is therefore PD stable. Its
ordinary algebra quotient is the tensor
product, and its universal property is the
PD pushout universal property. Differential
stability follows from the same relations
for $dc$.

Finally tensor (AT.25) with $X$ over $E(M)$.
The signed formula $x\otimes e\mapsto
(-1)^{|x|}x\otimes H(e)$ is well defined by
the second identity in (AT.25). It contracts
the tensor product onto $X$, whose retraction
is $x\otimes e\mapsto x f(\pi(e))$.
This proves (AT.23) without a flatness
assumption. The argument can be applied
separately in each simplex, proving the
stated simplicial form of the assertion.
\end{proof}

\paragraph{Exact image lattice and local coefficient tails.}
The ordinary polynomial map
\[
 F(M)\longrightarrow
 \Lambda_R(M_x)\otimes_R\Gamma_R(M_y)
\]
is injective over $R$. For an ordered basis and a
monomial $x^\epsilon y^\beta$ of arity
$m=|\epsilon|+|\beta|$, it sends
$y^\beta$ to
$(\prod_i\beta_i!)\prod_i\gamma_{\beta_i}(y_i)$.
Consequently the exact coefficient lattice of
(AT.21) on that divided-power/exterior basis
monomial is
\[
 \left(\prod_i\beta_i!\right)D_m.
 \tag{AT.26}
\]
No saturation of this lattice is made.

Over $R=\mathbb Z_{(p)}$ and $N=p^e$, one has
$D_m=I_{mN}$. The inclusion $D_m\subset I_{mN}$
follows from weights. For the reverse inclusion,
write $a=lN+s$ with $l\geq m$ and $0\leq s<N$.
The coefficient of $T^{[a]}$ in
$\gamma_l(T^{[N]})T^{[s]}$ is a $p$-adic unit:
Legendre's factorial valuation, or the no-carry
base-$p$ digit calculation, gives
\[
 v_p\!\left(\frac{(lN+s)!}{l!(N!)^l s!}\right)=0.
\]
The displayed element belongs to $D_m$ by its
divided-power index $l\geq m$, and hence supplies
$T^{[a]}$ after division by a unit of $R$.
Thus (AT.26) becomes
$(\prod_i\beta_i!)I_{mN}$ in this local case.

In arity one (AT.22) sends the coefficient-tail
elements $a x(v)$ and $a y(v)$ to
$a\alpha(v)$ and $a\,d\alpha(v)$, for $a\in I_N$.
It does not include the unscaled generators
$x(v),y(v)$ of the free odd PD disk. In particular,
this theorem does not extend a prescribed map
from that unscaled disk. Neither does it prove
that the relative attachments (AT.23) fill an
arbitrary unscaled horizontal cycle. These
precise maps and the lattice (AT.26) are the
scope of the construction.
